\section{Dirac and Lorentz}

The discrete walk reproduces relativistic motion. The Dirac dispersion, a finite light cone, and rotational
invariance all come out of the knit. They are not assumed. They are \emph{measured} off a run, read from the time
series of the tones the same way a physicist reads a wave off an oscilloscope. Special relativity is not one of the
eight atoms, and it is not added on top of them. It is what the knit looks like from far away.

\begin{principle}[Relativity is emergent]
Special relativity is not a base law. It is what the knit looks like from far away.
\end{principle}

Two plain definitions fix the targets before the derivations begin. The \emph{Dirac equation} is the law for an
electron-like particle in the regime where both quantum mechanics and special relativity hold at once. Its
\emph{dispersion} is the relation between a wave's energy and its wavelength, the curve $E(k)$ that says how fast a
ripple of a given size carries energy. \emph{Lorentz invariance} is the central rule of special relativity. The
laws look the same to every observer moving at a steady speed, and there is one finite top speed shared by all of
them. A \emph{light cone} is the reach of that top speed, the set of places a signal sent from here can possibly
arrive. The claim of this section is that each of these comes out of the knit on the mesh, measured rather than
imposed.

\subsection{The light cone is finite and exact}

The knit is collide then stream. Each beat the tones in a dock interact by the fixed reversible table, and then
each tone moves one dock along its site direction. The streaming half is a clean shift on the mesh. A tone that
sits in a dock this beat sits in a neighbour dock the next beat, and never skips ahead.

So a pulse spreads exactly one dock per beat. That is a finite, frame-independent top speed, the same in every
direction and the same for every observer who counts beats and docks. Set it to one. Write $z = 1$ for the
dynamical exponent, the exponent that ties a length scale to a time scale. A massless front obeys $x \sim t$, the
hallmark of a wave with a fixed speed, not the $x \sim \sqrt{t}$ of diffusion.

\begin{result}[The discrete light cone]
The light cone is one dock per beat. For the massless case it is an exact permutation of docks, a shift on the
mesh. There is no limit and no approximation, and continuity is never invoked.
\end{result}

This is the cleanest result in the section. The massless light cone is \emph{literally} the discrete stream. It is
not the long-wavelength shadow of some continuous law that lives underneath. The streaming half of the knit is the
light cone, exactly, at every scale down to the single dock.

\subsection{The Dirac dispersion, measured}

To read the dispersion off the substrate, run the simplest non-trivial walk that still carries the Dirac structure.
Take a two-direction walk, a left-mover and a right-mover, each a tone streaming along its own site. Let the collide
half of the knit mix them at a fixed rate. Track one dock's tone over many beats, then Fourier transform the time
series. The peaks of that transform trace out the dispersion $E(k)$, energy against wavenumber, with nothing put in
by hand beyond the knit itself.

Two cases come out, depending on whether the collide half mixes the two movers at all.

\begin{table}[ht]
\centering
\begin{tabular}{@{}llll@{}}
\toprule
case & knit & measured dispersion & rest mass \\
\midrule
massless & pure stream & $E(k) = k$ exactly & $0$ \\
massive & stream plus a site rotation by angle $m$ & $\cos E = \cos m \, \cos k$ & $E(0) = m$ \\
\bottomrule
\end{tabular}
\caption{The measured dispersion of the two-direction walk. The massive case is the lattice Dirac dispersion.}
\label{tab:dirac-dispersion}
\end{table}

The massive case is exactly the lattice Dirac dispersion. The mass is not a separate energy bolted on. It is the
coarse-grained mixing rate between the two movers, the angle by which the collide half rotates the sites each beat.

\begin{principle}[Mass is a mixing rate]
Mass is how fast the left-mover and the right-mover trade places per beat.
\end{principle}

So the Dirac dispersion is not put in. It is the knit's own measured dispersion. The complex quantum amplitudes of
the usual Dirac formalism are the emergent description of the real walk's oscillation modes, the bookkeeping that
tracks how the two movers oscillate into each other. The base is the walk, and the amplitudes ride on top of it.

This two-direction walk is the one-dimensional, two-component demonstrator, the cleanest reading available. The full
twenty-four-site version is the bulk wave on the mesh, and its isotropy is established in the next subsections.

\subsection{The full 3+1D Dirac operator}

On the $\{3,4,3,4\}$ mesh the twenty-four sites of a dock split, under triality, into three octets, $8v + 8s + 8c$.
The vector octet $8v$ is the force sector. The two spinor octets $8s$ and $8c$ are the two chiralities of matter,
the left-handed and right-handed halves of an electron-like field. The Dirac operator acts on these two spinor
sectors.

The decisive structural fact is about the \emph{square} of the operator. At every momentum the squared Dirac
operator is the relativistic scalar, $H^2 = p^2 + m^2$. The Hamiltonian itself is a matrix that mixes the two
chiralities, but its square collapses to a single number per momentum, the relativistic energy relation with no
matrix structure left.

\begin{theorem}[The squared Dirac operator is the relativistic scalar]
The Dirac operator built from the twenty-four site directions of a dock satisfies $H^2 = p^2 + m^2$ at every
momentum. This is not assumed. It is what the sites' Clifford algebra delivers once the operator is assembled from
the twenty-four directions, and it is the hinge from which the light cone, the magnetic moment, and the mass all
hang.
\end{theorem}

That $H$ squares to a clean scalar is the emergent relativistic structure. It is a fact about how the twenty-four
directions multiply, the Clifford algebra of the site set, and not a postulate. Two further consequences ride on
this single fact, and each is separately measured. The magnetic moment of the electron-like field comes from the
zero mode of this squared operator placed in a field, the spin gap that a spinless scalar in the same field does
not have. The mass is the coupling between the $8s$ and $8c$ chiralities, read off this same squared operator as
the part of $H$ that fails to commute with the chirality grading. Both are taken up in their own sections, the
magnetic moment under the g-factor and the mass under mass and chirality.

\begin{principle}[The hinge]
$H^2 = p^2 + m^2$ is the hinge. The light cone, the g-factor, and the mass all hang on it.
\end{principle}

\subsection{Isotropy to order four}

A lattice has preferred directions, its axes. Yet space, as we measure it, looks round, with no axis singled out.
The twenty-four sites deliver that roundness far better than a cube does, and the reason is a counting fact about
the site set.

The wave dispersion on the mesh is a sum over the site directions. Writing $d$ for a site direction and $k$ for the
wavevector,
\[
\omega^2(k) = \sum_{d} \bigl( 1 - \cos(k \cdot d) \bigr).
\]
Expanding the cosine to fourth order, the term that could spoil roundness is the fourth moment of the direction set,
the sum of $d^4$ against the sum of $d^2 d^2$. For the $D_4$ sites these balance exactly. The sum of $d^4$ equals
three times the sum of $d^2 d^2$, and both equal twelve. So the dispersion is isotropic to order four. The
axis-versus-diagonal anisotropy, the difference between a wave running along an axis and one running along a
diagonal, is exactly zero.

\begin{table}[ht]
\centering
\begin{tabular}{@{}llll@{}}
\toprule
lattice & directions & order-4 anisotropy & light cone shape \\
\midrule
cubic / hypercubic & $6$ or $8$ & about $8$ to $9$ percent & faceted, ratio $0.50$ \\
$D_4$ sites & $24$ & zero & round, ratio $1.00$ \\
\bottomrule
\end{tabular}
\caption{Order-four anisotropy of the wave dispersion. The twenty-four sites give an exactly round light cone where
a cube gives a faceted cross.}
\label{tab:isotropy}
\end{table}

This was measured both ways, analytically by the moment identity and on the GPU by simulation, and the two agree.
The twenty-four-site light cone is round. The eight-direction cubic one is a faceted cross, with rays poking out
along the diagonals.

\subsection{Triality kills the cubic anisotropy}

Why is the fourth moment isotropic at all? The answer is triality once more. Triality is the extra outer symmetry of
the twenty-four-cell, the order-three permutation of the three octets, and it is tied to the exceptional group
$F_4$. The point is an invariant count. $F_4$ has only one independent degree-four invariant, where the hypercubic
group $B_4$ has two. With only one invariant there is no room for a quantity that distinguishes an axis from a
diagonal, because that distinction would need a second independent degree-four invariant to live in.

\begin{principle}[No room for a preferred axis]
Triality leaves no room for a cubic anisotropy at order four. One invariant cannot encode a preferred axis.
\end{principle}

So the discrete $F_4$ symmetry of the sites restores the continuous rotational symmetry, the $\mathrm{SO}(8)$
roundness, in the long-wavelength limit. Isotropy is emergent, and it is sharper here than on any cube. The discrete
symmetry is not merely good enough, it is better. This is the emergent-symmetry result. Triality kills the
degree-four anisotropy, and the site symmetry stays round all the way out to order six before any lattice artifact
can appear.

\subsection{Boosts and the cusp caveat}

A boost is the change of frame that mixes space and growth, the relativistic tilt that trades position for time. With
a finite invariant top speed, $z = 1$, and an isotropic light cone, the boost structure that the mesh supports in
the infrared is the Lorentz one. The roundness of the light cone and the single top speed are exactly the two
ingredients a boost needs.

One honest subtlety remains. The bulk isotropy established above is near-perfect, the roundness of the full
four-dimensional mesh. But physical matter does not ride on the full bulk. It rides on the cubic cusp, the
$\{4,3,4\}$ slice that is physical space. That slice is anisotropic at order four, with a tiny residual that scales
as the square of the energy in Planck units, $(E / E_{\mathrm{Planck}})^2$.

\begin{result}[Bulk isotropy versus cusp isotropy]
The cusp is anisotropic at order four, a tiny residual of size $(E / E_{\mathrm{Planck}})^2$. So matter-propagation
isotropy is the cubic one, round in the infrared with a vanishing ultraviolet anisotropy. It is distinct from the
bulk's near-perfect roundness.
\end{result}

This bulk-cusp split is the honest part of the result, and it is carried forward into the cosmology section, where
the tiny cusp residual becomes a concrete prediction rather than a caveat.

\subsection{Why masslessness needs opposite faces}

There is a clean geometric precondition for an exactly massless mode, and it is worth stating because a generic
lattice fails it. The dock must have opposite faces.

On an odd-faced dock the instruction to go straight lands between two faces, with no facet pointing the way. That
forces a small zigzag every beat, an irreducible scatter that mixes the movers whether they should mix or not. The
lightest excitation then carries a floor mass set by that forced scatter, and nothing on such a dock is exactly
massless.

The twenty-four-cell is even-faced. Its twenty-four facets form twelve opposite pairs. A signal can go dead
straight, from a facet to its antipode, with mixing rate exactly zero. So the $\{3,4,3,4\}$ mesh hosts the
exactly-massless $E = k$ mode that a generic odd tiling could not support.

\begin{principle}[Masslessness rests on antipodal facets]
The masslessness rests on the dock's antipodal structure. The geometry is not cosmetic. An odd-faced dock would
force a floor mass on every excitation.
\end{principle}

\subsection{One mechanism, three limits}

The telegraph picture ties the whole section together with a single knob, the scatter rate $\lambda$, the rate at
which the collide half mixes the two movers. A persistent walk that scatters at rate $\lambda$ is exactly the
telegrapher's equation, which in one dimension is the Dirac equation. Turning the one knob sweeps through three
distinct regimes.

\begin{table}[ht]
\centering
\begin{tabular}{@{}lll@{}}
\toprule
$\lambda$ & regime & scaling \\
\midrule
$0$ & massless wave, light cone & $z = 1$ \\
finite & relativistic massive Dirac & $z = 1$, mass proportional to $\lambda$ \\
infinite & diffusion & $z = 2$ \\
\bottomrule
\end{tabular}
\caption{One mechanism, three limits. The scatter rate $\lambda$ interpolates from a massless wave through the
massive Dirac walk to ordinary diffusion.}
\label{tab:telegraph}
\end{table}

At zero scatter the walk is a pure wave on its light cone. At finite scatter it is the relativistic massive Dirac
walk, with mass proportional to $\lambda$. At infinite scatter the persistence is washed out and the walk becomes
ordinary diffusion, with $z = 2$. One mechanism, three limits, and the physical regime sits at finite $\lambda$
where the Dirac walk lives.

\subsection{Why it matters}

This closes the bridge from the base to quantum mechanics. Relativistic quantum mechanics is an emergent limit of
the discrete substrate, not a separate postulate sitting beside it. The massless light cone is exactly discrete,
needing no continuity. Mass emerges as a mixing rate, the coarse-grained speed at which the two movers trade places.
And space looks round, more so than on any cube, because triality forbids a degree-four anisotropy.

\begin{principle}[The walk and the cone]
The light cone is how fast the fire spreads. The Dirac walk is the rhythm of the song, dock to dock.
\end{principle}

The spinor that propagates along this walk is the one the sites carry, established under spin and fermions. The two
consequences of the hinge, the magnetic moment and the mass, are taken up under the g-factor and under mass and
chirality. The tiny cusp residual returns under cosmology and hierarchy, where it becomes a measurable signature
rather than a footnote.
